\scanpage{14}
\mspage{11}

une structure de modèle de $1$-forme, \unsure{\{}plus précisément, de $1$-segment\unsure{\}}, avec $\delta L=\{x_0,x_1\}$.

Si on échange les rôles de $x_0,x_1$ (i.e. choisit une autre ``origine'' du modèle de segment), alors la relation d'ordre est remplacée par son opposée (NB les conditions de divisibilité et de bifiltration \inserted{locale} sont autoduales). Ainsi, une structure de \inserted{modèle de} segment équivaut à une structure de \uline{biordre} (paire d'ordres \unsure{opposés|opposée}) qui soit telle que pour chacun \inserted{(ou l'un)} des ordres qui la composent, il y ait un plus petit et plus grand élément distincts l'un de l'autre (\unsure{i.e.|si} $\operatorname{Card}L\unsure{>|\ne}1$ -- en fait, $\operatorname{card}L=\infty$), divisible, et \unsure{loc$^{\mathrm t}$} filtrante et cofiltrante.

\mspage{13}

\rotatebox{90}{\unsure{EXTRA STRONG}}
\rotatebox{90}{\unsure{EXTRA STRONG}}
